\documentclass[11pt]{article}

\usepackage{CJKutf8}
\usepackage{math_article}
\usepackage{series_ra}

\renewcommand{\ArticleNumeral}{Sup}
\renewcommand{\ArticleTitle}{一致可积性与 $L^1$ 收敛}
\renewcommand{\ArticleAuthor}{Anqiao Ouyang}
\renewcommand{\ArticleDate}{August 2026}

\begin{document}
\begin{CJK*}{UTF8}{gbsn}

\maketitle


\section*{引言}


在讨论函数列的各种收敛方式时，我们很容易产生一种期待：

如果 $f_n$ 已经在越来越强的意义下接近 $f$，那么积分是否也应该随之接近？

也就是说，我们什么时候能够从

\[
f_n\to f
\]

推出

\[
\int_X f_n\,d\mu
\longrightarrow
\int_X f\,d\mu?
\]

这个问题在 Lebesgue 积分理论里并不陌生。单调收敛定理通过非负性与单调性解决它；控制收敛定理则要求存在一个统一的可积控制函数 $g$，满足

\[
|f_n|\leq g.
\]

这些条件都在做同一件事情：阻止函数列的“质量”以某种方式逃走。

但是控制函数的要求有时过强。

一个函数列可能根本不存在共同的可积上界，却仍然具有足够好的积分行为。此时真正重要的对象不是某一个控制函数，而是整个函数族对“大函数值区域”的统一控制。

这就是\textbf{一致可积性}（uniform integrability）。


\section{集中的质量}


先看一个最简单的例子。

在 $([0,1],\mathcal B,m)$ 上定义

\[
f_n(x)
=
n\mathbf 1_{(0,1/n)}(x).
\]

对任意固定的 $x>0$，当 $n$ 足够大时，

\[
x\notin(0,1/n),
\]

因此

\[
f_n(x)\to0.
\]

所以

\[
f_n\to0
\qquad\text{a.e.}
\]

同时，对于任意 $\varepsilon>0$，当 $n>\varepsilon$ 时，

\[
\{|f_n|>\varepsilon\}
=
(0,1/n),
\]

从而

\[
m(\{|f_n|>\varepsilon\})
=
\frac1n\to0.
\]

因此它也依测度收敛到 $0$。

但是

\[
\int_0^1|f_n|\,dx = n\cdot\frac{1}{n} = 1.
\]

所以

\[
|f_n|_{L^1}=1
\]

始终不变，更不可能有

\[
f_n\to0
\quad\text{in }L^1.
\]

问题出在哪里？

并不是函数列的总积分越来越大。事实上它始终等于 $1$。

真正发生变化的是这部分积分被压缩到了越来越小的集合上：

\[
\text{高度 }n
\qquad
\times
\qquad
\text{宽度 }\frac1n.
\]

集合越来越小，函数值却越来越大，两种变化恰好抵消。

因此，单纯知道

\[
\sup_n|f_n|_1<\infty
\]

是不够的。

我们需要禁止这种“有限的质量集中在无限高的尖峰中”的现象。


\section{定义}


设 $(X,\mathcal F,\mu)$ 是测度空间，$\mathcal G\subset L^1(\mu)$ 是一族可积函数。

如果

\[
\lim_{K\to\infty}
\sup_{f\in\mathcal G}
\int_{\{|f|>K\}}|f|\,d\mu
=
0,
\]

那么称 $\mathcal G$ \textbf{一致可积}。

对于函数列 $\{f_n\}$，这就是

\[
\lim_{K\to\infty}
\sup_n
\int_{\{|f_n|>K\}}
|f_n|\,d\mu
=
0.

\]

这个定义关注的不是整个积分

\[
\int_X|f_n|\,d\mu,
\]

而是函数值非常大的部分：

\[
\int_{\{|f_n|>K\}}|f_n|\,d\mu.
\]

固定一个很大的 $K$ 后，每个函数都可能存在超过 $K$ 的部分。但如果函数族一致可积，那么只要 $K$ 足够大，这些部分所携带的积分质量就能够\textbf{同时}变得任意小。

这里最关键的是

\[
\sup_n.
\]

对于任意单个 $f_n\in L^1$，我们本来就有

\[
\int_{\{|f_n|>K\}}|f_n|\,d\mu\to0.
\]

一致可积性要求这个收敛对整个函数族一致成立。

因此它并不是“每个函数都可积”的另一种写法，而是一个真正的\textbf{函数族性质}。概率论中也通常以完全相同的尾积分形式定义 uniform integrability。


\section{有界性}


一致可积首先会推出

\[
\sup_{f\in\mathcal G}
|f|_1<\infty.
\]

证明并不复杂。

由一致可积性，可以取某个 $K>0$，使得

\[
\sup_{f\in\mathcal G}
\int_{\{|f|>K\}}|f|\,d\mu
<1.
\]

把积分分成两部分：

\[
\int_X|f|\,d\mu
=
\int_{\{|f|\leq K\}}|f|\,d\mu
+
\int_{\{|f|>K\}}|f|\,d\mu.
\]

如果 $\mu(X)<\infty$，则第一项满足

\[
\int_{\{|f|\leq K\}}|f|\,d\mu
\leq
K\mu(X).
\]

于是

\[
|f|_1
\leq
K\mu(X)+1.
\]

因此

\[
\sup_{f\in\mathcal G}|f|_1<\infty.
\]

但反方向不成立。

上一节的

\[
f_n=n\mathbf1_{(0,1/n)}
\]

已经给出了反例。虽然

\[
\sup_n|f_n|_1=1,
\]

但只要 $n>K$，

\[
\{|f_n|>K\}
=
(0,1/n),
\]

因此

\[
\int_{\{|f_n|>K\}}|f_n|\,dx
=
1.

\]

于是对任何 $K$，

\[
\sup_n
\int_{\{|f_n|>K\}}|f_n|\,dx
=
1.

\]

它并不一致可积。

所以

\[
\boxed{
\text{一致可积}
\Longrightarrow
L^1\text{ 一致有界}
}
\]

而逆命题不成立。

两者之间相差的，正是对\textbf{质量集中}的控制。


\section{小集合}


在有限测度空间中，一致可积性还有一种非常重要的理解方式。

如果 $\mathcal G$ 一致可积，那么对任意 $\varepsilon>0$，存在 $\delta>0$，使得对于所有满足

\[
\mu(E)<\delta
\]

的可测集 $E$，都有

\[
\sup_{f\in\mathcal G}
\int_E|f|\,d\mu
<
\varepsilon.
\]

也就是说：

\begin{quote}
一个足够小的集合，不可能承载函数族中任何一个函数的大量 $L^1$ 质量。
\end{quote}


证明可以直接从截断开始。

由一致可积性，取 $K$ 使得

\[
\sup_{f\in\mathcal G}
\int_{\{|f|>K\}}|f|\,d\mu
<
\frac{\varepsilon}{2}.
\]

对于任意可测集 $E$，

\[
\int_E|f|\,d\mu
=
\int_{E\cap\{|f|\leq K\}}|f|\,d\mu
+
\int_{E\cap\{|f|>K\}}|f|\,d\mu.
\]

第一项至多为

\[
K\mu(E),
\]

第二项至多为

\[
\frac{\varepsilon}{2}.
\]

所以只需要令

\[
\delta=\frac{\varepsilon}{2K},
\]

当 $\mu(E)<\delta$ 时便有

\[
\int_E|f|\,d\mu<\varepsilon.
\]

这其实比尾积分定义更接近一致可积性的几何意义。

普通的 $L^1$ 有界只控制

\[
\int_X|f|,
\]

而一致可积还要求：\textbf{积分质量不能任意集中到越来越小的集合上。}

在有限测度空间中，加上 $L^1$ 一致有界后，这个“小集合上的一致绝对连续性”也可以反过来刻画一致可积性。


\section{$L^p$ 控制}


一致可积本身看起来仍然有些抽象，但有一个非常常用的充分条件。

设

\[
1<p<\infty
\]

并且函数族 $\mathcal G$ 满足

\[
\sup_{f\in\mathcal G}|f|_p\leq C.
\]

那么 $\mathcal G$ 在 $L^1$ 意义下是一致可积的。

事实上，在集合

\[
A_K=\{|f|>K\}
\]

上有

\[
|f|
\leq
\frac\{|f|^p\}{K^{p-1}},
\]

因此

\[
\int_{A_K}|f|\,d\mu
\leq
\frac1{K^{p-1}}
\int_{A_K}|f|^p\,d\mu
\leq
\frac\{|f|_p^p\}{K^{p-1}}.
\]

所以

\[
\sup_{f\in\mathcal G}
\int_{\{|f|>K\}}|f|\,d\mu
\leq
\frac{C^p}{K^{p-1}}.
\]

因为 $p>1$，

\[
K^{1-p}\to0,
\]

于是得到一致可积。

这里也能看出为什么 $p=1$ 是一个特殊的临界位置。

若只有

\[
\sup_n|f_n|_1<\infty,
\]

右侧不会产生任何随 $K\to\infty$ 消失的因子。

而一旦多出任意一点高阶可积性

\[
p>1,
\]

就产生了

\[
K^{1-p},
\]

从而自动控制高函数值尾部。


\section{控制收敛}


控制收敛定理的假设也会自动给出一致可积性。

假设存在

\[
g\in L^1(\mu)
\]

使得对所有 $n$，

\[
|f_n|\leq g
\qquad\text{a.e.}
\]

如果

\[
|f_n|>K,
\]

那么必然也有

\[
g>K.
\]

因此

\[
\{|f_n|>K\}
\subseteq
\{g>K\},
\]

并且

\[
\int_{\{|f_n|>K\}}|f_n|\,d\mu
\leq
\int_{\{g>K\}}g\,d\mu.
\]

由于 $g\in L^1$，

\[
\int_{\{g>K\}}g\,d\mu\to0.
\]

所以

\[
\{f_n\}
\]

一致可积。

因此，控制收敛定理中的统一控制函数实际上提供了比证明积分收敛所必需的条件更强的结构。

它给出的是

\[
|f_n|\leq g\in L^1,
\]

而真正负责控制积分尾部的是

\[
\{f_n\}\text{ uniformly integrable}.
\]

这也是 Vitali 收敛定理能够比控制收敛定理更灵活的原因之一。统一可积性可看成可积控制条件的一种放松。


\section{Vitali 定理}


现在可以回到最初的问题。

设

\[
\mu(X)<\infty
\]

且

\[
f_n\to f
\qquad\text{in measure}.
\]

仅有依测度收敛并不能推出 $L^1$ 收敛，尖峰函数已经说明了这一点。

Vitali 收敛定理指出，缺失的条件正是一致可积性：

\[
\boxed{
f_n\to f\text{ in measure}
+
\{f_n\}\text{ uniformly integrable}
\Longrightarrow
f_n\to f\text{ in }L^1.
}
\]

事实上，在这个框架下还可以得到相应的逆向刻画：当 $f_n\to f$ 依测度时，$L^1$ 收敛与一致可积性本质上是等价的。概率测度版本的 Vitali 定理通常就以这种形式陈述。

为什么这个结论成立？

取 $\eta>0$，定义坏集

\[
A_n
=
\{|f_n-f|>\eta\}.
\]

由于

\[
f_n\to f
\qquad\text{in measure},
\]

有

\[
\mu(A_n)\to0.
\]

再把 $L^1$ 误差拆成

\[
\int_X|f_n-f|\,d\mu
=
\int_{X\setminus A_n}|f_n-f|\,d\mu
+
\int_{A_n}|f_n-f|\,d\mu.
\]

在 $X\setminus A_n$ 上，

\[
|f_n-f|\leq\eta,
\]

于是

\[
\int_{X\setminus A_n}|f_n-f|\,d\mu
\leq
\eta\mu(X).
\]

而 $A_n$ 的测度趋于零。

一致可积性恰好保证，在这样越来越小的集合上，

\[
\int_{A_n}|f_n|
\]

以及相应的极限函数部分都不能保持一个固定的正质量。

因此第二项也趋于零。

最后先令 $n\to\infty$，再令 $\eta\to0$，便得到

\[
|f_n-f|_1\to0.
\]

这个证明揭示了 Vitali 定理真正的结构：

\[
\boxed{
\begin{array}{c}
\text{依测度收敛}\
\Downarrow\
\text{坏集的测度趋于 }0
\end{array}
\qquad+\qquad
\begin{array}{c}
\text{一致可积}\
\Downarrow\
\text{小集合无法携带大量积分}
\end{array}
}
\]

两部分恰好拼成了 $L^1$ 收敛。


\section{两类控制}


从这里再看控制收敛定理，会发现两者的逻辑其实非常接近。

控制收敛使用一个固定函数

\[
g\in L^1
\]

同时控制所有 $f_n$。

这是一种\textbf{点态控制}：

\[
|f_n(x)|\leq g(x).
\]

一致可积性则不要求存在这样的 $g$。

它只要求从积分意义上统一控制函数族的尾部：

\[
\sup_n
\int_{\{|f_n|>K\}}|f_n|
\to0.
\]

这是\textbf{质量控制}。

因此可能存在一致可积的函数列，却不存在任何共同的 $L^1$ 控制函数。Vitali 定理仍然能够处理这种情形；这也是它严格超出控制收敛定理直接适用范围的地方。Stanford 的概率论讲义也给出了这类“一致可积但不存在共同可积控制变量”的例子。

所以，两种定理可以放在同一个逻辑框架中理解：

\[
\text{DCT}
:
\quad
\text{共同 }L^1\text{ 控制}
\Longrightarrow
\text{一致可积}
\Longrightarrow
L^1\text{ 极限可控},
\]

而 Vitali 直接从中间这一层开始。


\section{无限测度}


这里还有一个容易被忽略的边界。

前面的有限测度假设不是完全无关紧要的。

在

\[
(\mathbb R,\mathcal B,m)
\]

上考虑

\[
f_n(x)
=
\frac1n\mathbf1_{[0,n]}(x).
\]

显然

\[
|f_n|_1=1.
\]

并且由于

\[
|f_n|_\infty=\frac1n,
\]

实际上

\[
f_n\to0
\]

甚至是一致收敛的，所以当然也依测度收敛。

另一方面，当 $K>1$ 时，

\[
\{|f_n|>K\}=\varnothing
\]

对所有 $n$ 成立，因此按照尾积分定义，$\{f_n\}$ 是一致可积的。

但是

\[
|f_n|_1=1
\]

仍然没有趋于零。

这里的质量并没有集中到越来越小的集合，而是发生了另一种现象：

\[
\text{质量向越来越大的空间范围扩散。}
\]

在有限测度空间里，这种行为没有足够的空间发生；在无限测度空间中则需要额外控制。

因此，对一般测度空间而言，仅仅控制

\[
|f_n|\to\infty
\]

方向上的尾部还不够，还需要防止积分质量向空间的“远处”逃逸。

一般版本的 Vitali 型结论因此还会加入一种 tightness 条件：能够找到有限测度区域，使其外部的积分质量对所有 $n$ 都一致地小。一般测度空间版本确实需要在尾值控制之外增加这种空间方向的控制。

这两个现象最好区分：

\[
\boxed{
\begin{aligned}
\text{concentration}
&:\quad
\text{质量挤进越来越小的集合},\
\text{escape}
&:\quad
\text{质量跑向越来越大的空间区域}.
\end{aligned}
}
\]

一致可积性首先排除前者，而在无限测度空间上还必须额外处理后者。


\section*{总结}


一致可积性讨论的并不是单个函数是否可积，而是一个函数族的积分质量是否受到统一控制。

它要求

\[
\lim_{K\to\infty}
\sup_n
\int_{\{|f_n|>K\}}
|f_n|\,d\mu
=
0,
\]

也就是函数列不能通过越来越高、越来越窄的尖峰保存固定的积分质量。

因此，

\[
\sup_n|f_n|_1<\infty
\]

本身并不足够，而任意 $p>1$ 的统一 $L^p$ 控制都会自动给出一致可积性。

更重要的是，在有限测度空间中，Vitali 收敛定理把它与函数列的收敛行为联系起来：

\[
\boxed{
f_n\to f\text{ in measure}
\quad+\quad
\{f_n\}\text{ uniformly integrable}
\quad\Longrightarrow\quad
f_n\to f\text{ in }L^1.
}
\]

依测度收敛告诉我们，误差较大的点越来越少。

一致可积性告诉我们，这些越来越少的点不能仍然携带大量积分。

两个条件控制的是两个不同的问题，而它们组合起来，正好补上了从依测度收敛到 $L^1$ 收敛之间缺失的一步。

从这个角度看，一致可积性不是又一种新的“收敛方式”。

它更接近于一种\textbf{紧性条件}：它控制一个函数族的质量如何分布，并决定我们什么时候能够安全地把极限送进积分号。

\end{CJK*}
\end{document}
